% Par acción-reacción en el tira y afloja (DIAG-05: Tercera Ley de Newton).
% Figuras simples con trazos (sin fuentes/emoji) para evitar depender de
% paquetes extra en el build de CI (solo tikz estándar).
\documentclass[tikz,border=3pt]{standalone}
\begin{document}
\begin{tikzpicture}[>=Latex, thick, font=\large]
  \draw[line width=2.5pt, brown] (-2.2,0) -- (2.2,0);

  % Camila (izquierda)
  \draw (-3.2,0) circle (0.25);
  \draw (-3.2,-0.25) -- (-3.2,-1.1);
  \draw (-3.2,-0.4) -- (-2.6,-0.1);
  \draw (-3.2,-0.9) -- (-2.7,-1.4);
  \draw (-3.2,-0.9) -- (-3.7,-1.4);
  \node at (-3.2,0.6) {Camila};

  % Compañero (derecha)
  \draw (3.2,0) circle (0.25);
  \draw (3.2,-0.25) -- (3.2,-1.1);
  \draw (3.2,-0.4) -- (2.6,-0.1);
  \draw (3.2,-0.9) -- (2.7,-1.4);
  \draw (3.2,-0.9) -- (3.7,-1.4);
  \node at (3.2,0.6) {Compañero};

  % Par acción-reacción en el punto donde Camila sostiene la cuerda
  \draw[->, red!70!black] (-2.2,0) -- (-3.0,0) node[midway, above] {$80\,N$};
  \draw[->, blue!70!black] (-2.2,0) -- (-1.4,0) node[midway, above] {$80\,N$};
\end{tikzpicture}
\end{document}
